\section{The Compiled Iteration, Formally}
\label{sec:theory}

This section states the per-iteration dataflow as a composition of three indexed
operators and establishes four properties of it. The folded chance representation
computes the same values as an explicit chance chain. The branchless dual-lane update
computes \cref{eq:lanes} exactly. The depth schedule is the shortest among schedules
that execute edges in dependency-respecting groups. Unrescaled linear averaging has a
representability horizon that fixes the accumulator precision. Proofs are in
\cref{app:proofs}.

\paragraph{Three indexed operators.}
\label{sec:operators}

Fix a compiled game with node set $\nodes$, $|\nodes| = \numnodes$, and edge set
indexed $1, \dots, \numedges$. Write $p(e), c(e)$ for the parent and child of edge
$e$ and $q(e)$ for its slot. Let $d(h)$ be the depth of node $h$ and
\begin{equation}
    B_\ell = \{\, e : d(p(e)) = \ell \,\},
    \quad
    \mathcal{D} = \{\, \ell : B_\ell \neq \emptyset \,\},
    \quad
    D = |\mathcal{D}|,
    \label{eq:blocks}
\end{equation}
so the blocks $\{B_\ell\}_{\ell \in \mathcal{D}}$ partition the edges by parent depth.
Let $\strat_{\mathrm{ext}} \in \R^{\numslots + 1}$ be the strategy vector extended by a
sentinel coordinate pinned at $\strat_{\mathrm{ext}}(\bot) = 1$, and for a lane
$i \in \{1,2\}$ define the compiled slot map
\begin{equation}
    q_i(e) \;=\;
    \begin{cases}
        q(e) & \text{edge } e \text{ is a decision edge of player } i,\\
        \bot & \text{otherwise (opponent or chance edge)},
    \end{cases}
    \label{eq:slotmap}
\end{equation}
which is a build-time constant. Its lane-free analogue $q_{\mathrm{val}}(e)$
($q(e)$ on decision edges, $\bot$ on chance edges) drives the backward pass.
Three operators then generate the iteration. For a
block $B_\ell$, a lane $i$, and vectors $x \in \R^{\numnodes}$, $y \in \R^{\numslots+1}$:
\begin{align}
    \fwd_{\ell,i}(x, y)(h)
        &=
        \begin{cases}
            x(p(e))\, y(q_i(e)) & h = c(e),\; e \in B_\ell,\\
            x(h) & \text{otherwise},
        \end{cases}
        \label{eq:opfwd} \\[2pt]
    \bwd_{\ell}(x, y)(h)
        &= x(h) + \sum_{e \in B_\ell \,:\, p(e) = h} y(q_{\mathrm{val}}(e))\, x(c(e)),
        \label{eq:opbwd}
\end{align}
\begin{equation}
    \acc_{\ell}(r, x, \rho)(q)
        = r(q)
        + \sum_{\substack{e \in B_\ell \,:\, q(e) = q \\ e \text{ a decision edge}}}\!\!
        s(e)\,\rho_{-i(e)}(p(e))\,\big(x(c(e)) - x(p(e))\big).
        \label{eq:opacc}
\end{equation}
Here $i(e)$ is the decision-edge owner, $s(e)=+1$ for player 1 and $-1$
for player 2, and $\rho_{-i(e)}$ selects the opponent's reach lane. The vector
$x$ holds chance-folded player-1 values, so no further chance factor is used.
Each operator uses gathers, elementwise arithmetic, and a scatter over $|B_\ell|$
entries, which is how \cref{alg:iteration} realizes them: \cref{eq:opfwd} as an
indexed assignment, \cref{eq:opbwd,eq:opacc} as \texttt{index\_add}. The forward pass is
$\fwd$ composed over $\ell \in \mathcal{D}$ in increasing depth, the backward pass is
$\bwd$ and $\acc$ composed in decreasing depth, and \cref{eq:opfwd} runs on the doubled
index range so both lanes advance in one call.

The forward pass assigns in \cref{eq:opfwd} where the backward pass accumulates.
Assignment is what frees it of buffer clearing. In a tree each non-root node is the
child of exactly one edge, so within a block every destination is written exactly
once and stale contents are unreachable. The backward operators accumulate instead, because a parent has several
children per block.

\paragraph{What the operators compute.}

\begin{proposition}[Folded chance is exact]
\label{prop:chance}
Initialize $\cfv \gets \cfv_{\mathrm{tmpl}}$ as in \cref{eq:template} and apply
$\bwd_\ell$ for $\ell \in \mathcal{D}$ in decreasing order, with every chance edge
assigned multiplier $1$ through $q_{\mathrm{val}}$. Then on termination, for every node
$h$,
\begin{equation}
    \cfv(h) \;=\; \sum_{z \succeq h,\; z \in \termnodes}
        \pi_c(h \!\to\! z)\, \pi_c(h)\, \reach^{\strat}_{1,2}(h \!\to\! z)\, u(z),
    \label{eq:chanceexact}
\end{equation}
that is, the chance-weighted continuation value at $h$. In particular at the root
$\cfv$ equals the expected payoff under $\strat$, and each terminal's chance product
$\pi_c(z)$ enters exactly once.
\end{proposition}

\Cref{prop:chance} says that folding saves work and leaves the computed values
unchanged. The chance chain is deleted from the dynamic pass and reappears only as
the build-time constant multiplying $u(z)$ in \cref{eq:template}.

\begin{proposition}[The dual-lane update is exact and branch-free]
\label{prop:lanes}
For every lane $i$, block $\ell$, and edge $e \in B_\ell$, the operator $\fwd_{\ell,i}$
of \cref{eq:opfwd} evaluated with the compiled slot map \cref{eq:slotmap} and
$y = \strat_{\mathrm{ext}}$ satisfies the case definition of \cref{eq:lanes} at
$c(e)$. The slot map is resolved at build time, so the per-iteration cost of the case
distinction is zero and no per-edge predicate is evaluated during execution.
\end{proposition}

\begin{proposition}[The depth schedule is shortest]
\label{prop:depth}
Call a schedule \emph{dependency-respecting} if it partitions the edges into ordered
groups such that no group contains two edges $e, e'$ with $c(e) = p(e')$, and $e$
precedes $e'$ whenever $c(e) = p(e')$. Then:
\begin{enumerate}[label=(\roman*),topsep=1pt,itemsep=0pt]
    \item the partition \cref{eq:blocks} is dependency-respecting and has $D$ groups;
    \item every dependency-respecting schedule has at least $D$ groups;
    \item the number of operator invocations per iteration is $c_1 + c_2 D$ for
    constants $c_1, c_2$ independent of $\numnodes$ and $\numedges$.
\end{enumerate}
\end{proposition}

Part~(iii) is the reason iteration time is nearly flat in tree size across four
decades (\cref{tab:matrix}). Growth enters the batch \emph{widths}, which the device
absorbs in parallel, while the number of dispatches stays fixed. Measured on the
suite, $c_1 \approx 31$ (the itemized dispatches of \cref{tab:opsteps} plus framework
bookkeeping) and $c_2 = 8$. The realized schedule applies $\acc$ once, flat over all
decision edges (\cref{sec:dataflow}), so only $\fwd$ and $\bwd$ contribute to $c_2$,
at four operations per level each. \Cref{tab:suite} confirms that $D$ matches the
tree depth on every game. The reference implementation's per-(stage, depth) partition
needs up to $100$ groups where \cref{eq:blocks} needs $8$.

Kim (2026) \cite{kim2026parallelizing} organizes its sequence-form iteration by
level as well, so its critical path is likewise proportional to depth. The
difference the timing measures is per-step work. There, one level applies a masked
sparse product against a precomputed level matrix, whose cost is set by the sparsity
pattern and the sparse kernel. Here, one level applies
\cref{eq:opfwd,eq:opbwd,eq:opacc} on dense contiguous ranges with build-time
indices. \cref{app:extra} localizes the resulting gap
to the sparse product and the index assignment inside its per-iteration
\texttt{normalize}.

\paragraph{The averaging horizon.}

Linear averaging \cite{tammelin2014solving} accumulates
$\bar{s} \mathrel{+}= t\,\reach_{\mathrm{own}}\,\strat$ on iteration $t$ and never
rescales. The accumulator therefore grows quadratically while each contribution stays
linear, and in finite precision the two rates collide.

\begin{proposition}[Representability horizon of unrescaled linear averaging]
\label{prop:horizon}
Let $S_T = \sum_{t=1}^{T} t\,x_t$ be the ideal real sum, with
$x_t\in[0,1]$, $S_T = \Theta(T^2)$, and $x_T\ge c>0$.
For a binary format with $p$-bit significand, while $S_T$ is in the normal range,
\begin{equation}
    \frac{T x_T}{\ulp(S_T)} = \Theta\!\left(\frac{2^p}{T}\right).
    \label{eq:horizon}
\end{equation}
Here $\ulp$ is the spacing in the binade containing the ideal sum (the
actual rounded accumulator may sit in a neighboring binade). Without $x_T\ge c$, only the $O(2^p/T)$ bound
follows. Thus $2^p$ is an order-of-magnitude resolution scale and should be read as such;
the exact stalling point depends on the rounded trajectory. A nonnegative increment is locally absorbed under
round-to-nearest if it is strictly below half the upward spacing
at the actual old accumulator; ties depend on the rounding rule.
\end{proposition}

GPU-CFR holds $\bar{s}$ in float64 for all update rules, while the other buffers
and iteration arithmetic use the configured dtype. The scales
$2^{24}\approx1.7\times10^7$ and $2^{53}\approx9\times10^{15}$ show the precision
difference. The cost is one wider slot vector with no change to the operator
sequence. This accumulator issue is separate from the graph-replay counter of
\cref{sec:cudagraph}. That float32 counter represents consecutive integers exactly
through $2^{24}$, and the host refills its starting index before each replay batch.
